\scanpage{74}
\mspage{11}
d) Les séries convergentes
\[
f=\sum_{i\geq0}a_i z^i\qquad(a_i\in\mathbb C)
\]
de rayon de convergence $\geq1$ (donc
définissant $\mathbb D^\circ\mathrel{\unsure{-|\to}}\mathbb C$) \inserted{avec $|f(z)|\leq1$ si $z\in\mathbb D$,} telles que la $f$
se prolonge par continuité en une
fonction $\mathbb D\mathrel{-}\deleted{\mathbb C}\mathbb D$, celle-ci ayant les
propriétés 1) 2) 3) de B) ci-dessus.

\uline{13.1.}\qquad\textcircled{6}

Soit $(X,C)$ un objet du groupoïde
de Teichmüller analytique complexe
à bord, avec la condition différentiable.
Pour fixer les idées, supposons $X$ connexe,
de type $g,\nu$. Considérons le découpé $X'$
de $X$ suivant $C$, bouchons ses trous \inserted{canoniquement en utilisant la structure torique de $\partial X'$, $*$}
% slanted
\margin{$*$ via la métrique riemannienne can. de $X'$}
d'où surface holomorphe compacte \inserted{sans bord $X''$}, à disques
\inserted{pointés} marqués, avec des \deleted{iso} \inserted{antiiso} de \deleted{ces} \inserted{certains de ces} disques
(\deleted{deux à deux} qui se groupent par paires, à l'exception
des $\nu$ disques correspondant à $\partial X$).

On peut considérer $X''$ comme le
normalisé d'une courbe \deleted{\unsure{alg}} \inserted{$\mathbb C$-alg $X'''$} ``stable''
(au sens de Mumford-Deligne), en
% archive sheet number: 76
